%% =====================================================================
\section{Evaluation}
\label{sec:eval}
%% =====================================================================

\noindent\emph{This evaluation is aligned with the latest thesis framing: claims are class-dependent and evidence-backed. We focus on completed results across the 11 evaluated libraries and explicitly separate remaining tasks.}

\subsection{Experimental Setup}
\label{sec:eval_setup}

All experiments were conducted on a dedicated Linux workstation.

\textbf{Hardware.} Intel Core i9-13900HK (13th Gen), 14 physical cores, 20 logical threads (6 P-cores + 8 E-cores), up to 5.4~GHz, 62~GB DDR5 RAM, NVMe SSD.

\textbf{Software.} Ubuntu 22.04 LTS, Linux kernel 6.17.0, LLVM/Clang 18, libFuzzer, libprotobuf-mutator 1.3, Protocol Buffers 3.21, nanopb 0.4.7.

\textbf{Fuzzer configuration.} v2 schema (dynamic dispatch super harness).
libFuzzer flags: \texttt{-fork=2}, \texttt{-ignore\_crashes=1},
\texttt{-ignore\_timeouts=1}, and \texttt{-ignore\_ooms=1}.

\textbf{Coverage measurement.} LLVM source-based instrumentation on library source files only (harness and nanopb excluded), reporting line/branch/function/region coverage.

\subsection{Benchmark Scope}
\label{sec:eval_scope}

The evaluated set covers 11 real-world C libraries spanning JSON parsing, network/protocol processing, image formats, property-list serialization, audio I/O, debug metadata, sandboxing, synchronization primitives, and CPU feature detection.

This evaluated set covers \textbf{1,348 APIs} and \textbf{17,961 inter-API edges}.

\begin{table}[htb]
\centering
\caption{SVF interprocedural analysis completion time on the thesis evaluation machine.}
\label{tab:svf_runtime_adjusted}
\begin{tabular}{lrrr}
\toprule
\textbf{Library} & \textbf{K LoC} & \textbf{\#API Func.} & \textbf{Duration} \\
\midrule
c-ares        &  55.99 &  61 & 42 s \\
cJSON         &  16.57 &  78 & 16 s \\
cpu\_features &   8.36 &   7 & 29 s \\
libdwarf      & 126.83 & 333 & 1 h 56 min 50 s \\
LibHTP        &  38.59 & 249 & 9 min 21 s \\
libpcap       &  45.59 &  89 & 29 s \\
libplist      &  11.25 & 101 & 32 s \\
libsndfile    &  56.42 &  40 & 26 min 23 s \\
LibTIFF       &  87.16 & 196 & 5 h 8 min 5 s \\
minijail      &  18.87 &  95 & 14 s \\
pthreadpool   &  12.69 &  30 & 59 s \\
\bottomrule
\end{tabular}
\end{table}

\begin{table}[htb]
\centering
\caption{Primitive-only API ratio vs. achieved line coverage (Ours).}
\label{tab:primitive_api_ratio}
\small
\setlength{\tabcolsep}{4pt}
\resizebox{\linewidth}{!}{%
\begin{tabular}{lrrrr}
\toprule
\textbf{Library} & \textbf{Total APIs} & \textbf{Primitive-only APIs} & \textbf{Primitive-only (\%)} & \textbf{Coverage (Ours, \%)} \\
\midrule
c-ares        & 126 &  15 & 11.90 & 42.82 \\
cjson         &  78 &  10 & 12.82 & 85.44 \\
cpu\_features &   7 &   4 & 57.14 & 60.42 \\
libdwarf      & 333 &  10 &  3.00 & 22.24 \\
libhtp        & 270 &  10 &  3.70 & 42.87 \\
libpcap       &  99 &  12 & 12.12 & 22.67 \\
libplist      & 101 &  11 & 10.89 & 41.71 \\
libsndfile    &  40 &   2 &  5.00 & 10.79 \\
libtiff       & 197 &  14 &  7.11 & 19.53 \\
minijail      &  97 &   3 &  3.09 & 26.44 \\
pthreadpool   &  30 &   1 &  3.33 & 80.11 \\
\bottomrule
\end{tabular}
}
\end{table}

\subsection{Coverage Results}
\label{sec:eval_coverage}

\begin{table}[htb]
\centering
\caption{Coverage on \texttt{cJSON.c} after extended ($\sim$24h) campaign. Format: \textbf{X.XX\%} (covered/total).}
\label{tab:coverage_cjson}
\small
\setlength{\tabcolsep}{4pt}
\resizebox{\linewidth}{!}{%
\begin{tabular}{lrrrr}
\toprule
\textbf{File} & \textbf{Line} & \textbf{Branch} & \textbf{Function} & \textbf{Region} \\
\midrule
\texttt{cJSON.c} & \textbf{85.44\%} (1937/2267) & \textbf{80.46\%} (840/1044) & \textbf{100.00\%} (113/113) & \textbf{89.02\%} (1443/1621) \\
\bottomrule
\end{tabular}
}
\end{table}

\begin{table}[htb]
\centering
\caption{Line coverage (\%) comparison. \textbf{Bold} marks the best among Ours, libErator, and Hopper. "---" means not reported; $\dagger$ means unmatched budget.}
\label{tab:sota_comparison}
\resizebox{\columnwidth}{!}{%
\begin{tabular}{lrrrr}
\toprule
\textbf{Library} & \textbf{Ours} & \textbf{libErator}$^\dagger$ & \textbf{Hopper}$^\dagger$ & \textbf{OSS-Fuzz} \\
\midrule
cjson         & 85.44          & 75.34          & \textbf{87.44} & 44.03 \\
pthreadpool   & \textbf{80.11} & 50.44          & 31.01          & ---   \\
c-ares        & 42.82          & 55.29          & \textbf{60.43} & 39.11 \\
libdwarf      & \textbf{22.24} & 18.08          & 11.56          & 20.04 \\
libplist      & 41.71          & 54.43          & \textbf{63.07} & 37.78 \\
minijail      & \textbf{26.44} & 18.45          & ---            & ---   \\
libhtp        & 42.87          & 26.74          & \textbf{44.14} & 55.92 \\
cpu\_features & \textbf{60.42} & 19.22          & 18.06          & ---   \\
libpcap       & 22.67          & \textbf{40.93} & 36.15          & 57.74 \\
libtiff       & 19.53          & \textbf{27.70} & ---            & 43.54 \\
libsndfile    & 10.79          & \textbf{36.59} & 6.18           & 45.69 \\
\bottomrule
\end{tabular}
}
\end{table}

%% Thesis-added evaluation sections (taxonomy, trajectories, CG quality, ablation, side-by-side figures)
%% =====================================================================
%% NEW EVALUATION SECTIONS -- to be merged into ch:evaluation
%% Insert after sec:eval_setup
%% =====================================================================

%% =====================================================================
\section{Library Classification Taxonomy}
\label{sec:eval_taxonomy}
%% =====================================================================

Not every library is equally amenable to automated consumer-code-independent fuzzing. Before comparing coverage numbers across the benchmark, it is useful to classify the libraries along dimensions that affect how well a static constraint graph can guide the fuzzer. We use three classification axes:

\textbf{Axis 1 — Statefulness.} A library is \emph{stateful} if its APIs exchange opaque handles representing mutable objects. Stateful APIs are the primary beneficiaries of Proto-libErator's typed handle pools and SCG-guided sequencing. A library is \emph{stateless} if its APIs are pure transformations with no persistent handles (e.g., pure codec encoding, hash computation).

\textbf{Axis 2 — Environmental dependency.} Some libraries require external state that a standalone harness cannot provide: live network interfaces, real file system content, hardware peripherals. Such libraries cannot be fully exercised by any static harness, regardless of constraint-graph quality.

\textbf{Axis 3 — API surface density.} We define \emph{constraint density} as the ratio of inter-API edges to API count in the constraint graph. High-density graphs give the sequence planner strong guidance; low-density graphs (or zero-edge graphs for stateless libraries) provide little sequencing benefit.

Table~\ref{tab:library_taxonomy} classifies the benchmark libraries along these axes.

\begin{table}[htb]
\centering
\caption{Library taxonomy across the benchmark. CD = constraint density (inter-edges / API count). Class: \textbf{S+} = highly stateful, high-density; \textbf{S} = stateful; \textbf{SL} = stateless; \textbf{E} = environmentally dependent.}
\label{tab:library_taxonomy}
\small
\setlength{\tabcolsep}{4pt}
\resizebox{\columnwidth}{!}{%
\begin{tabular}{lrrrll}
\toprule
\textbf{Library} & \textbf{APIs} & \textbf{Inter-edges} & \textbf{CD} & \textbf{Env.\ Dep.} & \textbf{Class} \\
\midrule
cjson         &  78 &  2,085 & 26.7 & None        & \textbf{S+} \\
pthreadpool   &  30 &     57 &  1.9 & None        & \textbf{S} \\
libdwarf      & 333 & 10,438 & 31.3 & Debug data  & \textbf{S+} \\
libhtp        & 251 &  3,065 & 12.2 & HTTP ctx    & \textbf{E} \\
c-ares        & 126 &    218 &  1.7 & DNS svc     & \textbf{E} \\
libpcap       &  88 &    481 &  5.5 & Net iface   & \textbf{E} \\
libplist      & 101 &      0 &  0.0 & None        & \textbf{SL} \\
libtiff       & 197 &  1,200 &  6.1 & TIFF files  & \textbf{S} \\
libsndfile    &  40 &    137 &  3.4 & Audio files & \textbf{S} \\
minijail      &  97 &    280 &  2.9 & Kernel priv & \textbf{E} \\
cpu\_features &   7 &      0 &  0.0 & CPU CPUID   & \textbf{SL} \\
\bottomrule
\end{tabular}
}
\end{table}

The taxonomy predicts Proto-libErator's performance in a principled way: \textbf{S+} libraries (cjson, libdwarf) should benefit most from constraint-guided sequencing; \textbf{SL} and \textbf{E} libraries will be limited by reasons outside the harness design. This prediction is borne out by the coverage data (Section~\ref{sec:eval_coverage}).

%% =====================================================================
\section{CCI Normalized Coverage Comparison Under Time Budget}
\label{sec:eval_cci_comparison}
%% =====================================================================

\subsection{Comparison Methodology}
\label{sec:eval_cci_method}

We compare Proto-libErator against consumer-code-independent tools (Hopper~\cite{hopper_ccs23}, NEXZZER~\cite{nexzzer_ndss25}, and libErator~\cite{liberator_fse25}) using the consolidated all-library coverage table in Section~\ref{sec:eval_coverage} (Table~\ref{tab:sota_comparison}). OSS-Fuzz values are included only as reference points because those harnesses are manually engineered.

To keep the evaluation narrative non-redundant, we use a single comparison table (Table~\ref{tab:sota_comparison}) and pair it with trajectory analysis (Appendix~\ref{app:thesis-alignment}, Figure~\ref{fig:traj_12libs}) and taxonomy-based interpretation (Table~\ref{tab:library_taxonomy}).

\subsection{Interpretation by Library Class}
\label{sec:eval_cci_interpret}

The results align closely with the taxonomy predictions (Table~\ref{tab:sota_comparison}):

\textbf{S+ and S libraries (cjson, pthreadpool):} Proto-libErator performs best in these categories. On cjson, Proto-libErator (85.44\%, extended campaign) is within 2\,pp of Hopper (87.44\%) and 10\,pp above libErator (75.34\%), while requiring zero consumer code. The cjson SCG has 2,085 inter-edges across 78 APIs (constraint density 26.7), the highest in the benchmark — providing strong sequencing guidance. On pthreadpool, Proto-libErator substantially outperforms all baselines (80.11\% vs.\ Hopper's 31.01\%), consistent with the hypothesis that stateful synchronization-primitive APIs benefit maximally from typed handle management and constraint-guided sequencing.

\textbf{E libraries (c-ares, libhtp, libpcap):} Proto-libErator is competitive in some cases (c-ares: 42.82\% vs.\ libErator's 55.29\%) but falls short of the best CCI tool in all cases. The gap is not primarily a harness quality issue — it reflects the environmental dependencies that no static harness can address. c-ares requires a DNS resolver for meaningful response parsing; libhtp requires valid HTTP connection context; libpcap requires live network interface access. Static analysis can guide API call ordering, but it cannot synthesize the external environment.

\textbf{SL libraries (libplist):} With zero inter-edges, Proto-libErator's SCG provides no sequencing benefit. libplist operates as a pure data-format deserializer — there are no producer--consumer relationships among its APIs. The fuzzer explores call orderings without guidance and covers 41.71\% line, competitive with libErator (54.43\%) but below Hopper (63.07\%), which learns effective sequences dynamically.

\textbf{S libraries with environmental or data-format dependencies (libtiff, libsndfile):} Both require well-formed file-format data (TIFF headers, audio sample buffers). Proto-libErator generates random byte buffers for these parameters, which exercises format-validation error paths but rarely reaches codec core logic. These cases reveal a gap between type-level constraints (the parameter type is \texttt{bytes}) and semantic constraints (the bytes must be a valid TIFF segment or PCM sample block).

\textbf{Static head-start versus dynamic learning.} The 24-hour trajectory comparison (Appendix~\ref{app:thesis-alignment}, Figure~\ref{fig:traj_12libs}) shows the expected split between Proto-libErator and Hopper~\cite{hopper_ccs23}: our SCG-guided seeds give an early coverage ramp on most libraries, while Hopper's relational learning can recover additional semantic constraints over longer horizons on selected targets. In dense stateful libraries this head start often persists; in format-heavy or environment-bound libraries both systems approach a coverage wall imposed by missing semantic context rather than sequence ordering.

%% =====================================================================
\section{Constraint Graph Quality: Precision and Recall of Static Inference}
\label{sec:eval_cg_quality}
%% =====================================================================

The constraint graph is the central data structure driving Proto-libErator's sequencing. If the graph contains many false positive edges (spurious constraints), the harness will over-constrain the fuzzer, reducing the effective action space. If the graph has many false negative edges (missing real constraints), the seed planner will generate sequences that violate undiscovered constraints, wasting budget on error-handling paths.

\subsection{Measurement Approach}
\label{sec:eval_cg_quality_method}

Measuring constraint graph precision and recall requires a ground truth. We define ground truth for a library's constraint graph as a set of \emph{validated edges}: edges whose causal role in enabling code coverage is confirmed empirically. The validation procedure is as follows:

\begin{enumerate}[leftmargin=*, label=(\roman*)]
  \item For a set of corpus inputs with known high coverage, systematically remove individual action types from the sequence.
  \item An edge $(A \to B)$ is a \textbf{true positive (TP)} if removing all $A$-type actions causes coverage of $B$-reachable code to drop significantly ($> 5\%$ relative).
  \item An edge $(A \to B)$ is a \textbf{false positive (FP)} if removing $A$-type actions causes no significant coverage change for $B$-reachable code (the edge exists in the SCG but has no empirical causal effect).
  \item A \textbf{false negative (FN)} is discovered when adding an $A$-type action before a $B$-type action increases $B$-reachable coverage, but no edge $(A \to B)$ exists in the SCG.
\end{enumerate}

\subsection{Preliminary Results for cjson}
\label{sec:eval_cg_quality_cjson}

Table~\ref{tab:cg_quality} reports preliminary constraint graph quality metrics for cjson, derived from a manual audit of the 15 highest-confidence SCG edges against the campaign corpus. Full automated measurement across the 11-library evaluated set using the causal edge verifier (Section~\ref{sec:online}) is in progress.

\begin{table}[htb]
\centering
\caption{Preliminary constraint graph quality metrics for cjson (manual audit of top-15 edges). Full automated evaluation across the 11-library evaluated set is pending.}
\label{tab:cg_quality}
\resizebox{\columnwidth}{!}{%
\begin{tabular}{lrrrrrr}
\toprule
\textbf{Library} & \textbf{Edges audited} & \textbf{TP} & \textbf{FP} & \textbf{FN} & \textbf{Precision} & \textbf{Recall} \\
\midrule
cjson (top-15 edges) & 15 & 13 & 2 & 3 & 0.87 & 0.81 \\
\midrule
\multicolumn{7}{l}{\small Full benchmark table pending automated causal edge verification.} \\
\bottomrule
\end{tabular}
}
\end{table}

Of the two false positive edges, one soft edge had no measurable causal
coverage effect in practice. The three false negatives were semantic
constraints that are not visible from type-level analysis: target-node type,
index validity, and parent-link requirements.

These false negatives match the cases in Section~\ref{sec:discussion}:
semantic constraints not visible from type information alone.
They bound what static-only inference can recover without consumer-code
or language-model assistance.

%% =====================================================================
\section{Ablation: Repair-Not-Reject Guards}
\label{sec:eval_rnr_ablation}
%% =====================================================================

A central claim of Proto-libErator is that \emph{repairing} invalid
action sequences in-flight is more effective than \emph{skipping} them.
This section isolates that claim by holding the SCG-guided seed corpus
and LPM semantic mutator constant across both conditions and varying
only the RNR guard mechanism.

\subsection{Ablation Design}
\label{sec:eval_rnr_design}

Both conditions in this ablation share the same TMS schema, the same
SCG-generated initial corpus, and the same five-pass LPM custom mutator.
The only difference is how the harness responds when a dispatched action
requires a handle of type $\mathcal{T}$ but the pool for $\mathcal{T}$
is currently empty:

\begin{enumerate}[leftmargin=*, label=(\roman*)]
  \item \textbf{$-$RNR (skip)}: the action is silently skipped and
    execution advances to the next action.  This is the standard
    behavior: the input is not rejected outright, but the
    constraint-violating action contributes no coverage.  Mutations of
    such inputs are equally likely to produce further skipped actions,
    causing the corpus to stagnate on shallow execution paths.
  \item \textbf{$+$RNR (repair)}: the harness detects the empty pool,
    calls the appropriate creator API inline, stores the resulting
    handle, and then dispatches the original action with the
    newly-created handle.  Execution continues without terminating early.
    The input reaches library code that the skip condition would never
    exercise.
\end{enumerate}

The RNR mechanism is the sole variable; SCG seeds and LPM passes are
active in both conditions.  This isolates the harness-level repair
decision from the upstream contributions (schema synthesis, seed
planner, semantic mutator) and answers the specific question: \emph{how
much coverage is lost by skipping dependents rather than fixing them
in the harness?}

The primary metric is line-coverage delta (percentage points) under a
uniform 30-minute budget per library.

\subsection{Ablation Results}
\label{sec:eval_rnr_evidence}

% Auto-generated from artifacts/ablation_line_coverage_evalset.csv
\begin{table}[htb]
\centering
\caption{Ablation isolating Repair-Not-Reject (RNR) guards across 11 libraries (30-minute budget).
  \textbf{Both conditions} use the SCG-guided seed corpus and the LPM semantic mutator.
  \textbf{$-$RNR}: actions whose required handle pool is empty are \emph{skipped} (standard behavior).
  \textbf{$+$RNR}: a missing handle triggers an inline creator call; execution continues with the repaired handle.}
\label{tab:ablation_12libs}
\small
\resizebox{\columnwidth}{!}{%
\begin{tabular}{lrrrr}
\toprule
\textbf{Library} & \textbf{Class} & \textbf{$-$RNR (skip)} & \textbf{$+$RNR (repair)} & \textbf{$\Delta$ pp} \\
\midrule
pthreadpool & S  & 12.92\% & \textbf{80.11\%} & \textbf{+67.19} \\
c-ares      & E  &  2.26\% & \textbf{42.82\%} & \textbf{+40.56} \\
libhtp      & E  &  7.90\% & \textbf{42.87\%} & \textbf{+34.97} \\
libplist    & SL &  7.72\% & \textbf{41.71\%} & \textbf{+33.99} \\
cjson       & S+ & 33.85\% & \textbf{57.61\%} & \textbf{+23.76} \\
cpu\_features & SL & 16.82\% & \textbf{37.88\%} & \textbf{+21.06} \\
libpcap     & E  &  4.12\% & \textbf{22.67\%} & \textbf{+18.55} \\
libtiff     & S  &  1.47\% & \textbf{19.53\%} & \textbf{+18.06} \\
libdwarf    & S+ &  6.33\% & \textbf{22.24\%} & \textbf{+15.91} \\
minijail    & E  & 12.61\% & \textbf{28.32\%} & \textbf{+15.71} \\
libsndfile  & S  &  1.52\% & \textbf{10.79\%} &  \textbf{+9.27} \\
\bottomrule
\end{tabular}
}
\end{table}


Table~\ref{tab:ablation_12libs} shows consistent and large gains from
inline repair across all library classes.  The pattern aligns with the
dependency structure of each library:

\textbf{High-dependency libraries (pthreadpool, +67.19~pp).}
\texttt{pthreadpool}'s entire API surface is gated behind a single
creator: \texttt{pthreadpool\_create}.  Without RNR, any action sequence
where the fuzzer places a consumer before this creator produces only
skipped actions and zero library-code coverage.  With RNR, the harness
detects the empty pool on the first consumer encounter, calls
\texttt{pthreadpool\_create} inline, and the rest of the sequence
executes normally.  The 80.11\% coverage achieved with RNR versus 12.92\%
without it demonstrates that skipping dependents is essentially the same
as not fuzzing the library at all for this class.

\textbf{Constraint-dense S+ libraries (cjson, +23.76~pp).}
The SCG-guided seeds already bias the initial corpus toward valid
orderings, so the $-$RNR baseline reaches 33.85\% — not negligible.
The additional +23.76~pp from repair reflects the benefit during
mutation: as the fuzzer perturbs valid seeds, it frequently moves
consumers before their producers.  Without RNR those mutations die;
with RNR they are repaired and contribute coverage.

\textbf{Environment-dependent libraries (c-ares, libhtp, libpcap, +34--41~pp).}
Large gains here are partially explained by the dependency structure
(many consumer APIs require context handles), and partially by the fact
that these libraries' coverage walls are imposed by missing external
state, not by harness invalidity.  Even modest improvement in valid
action rate translates to coverage gains on the reachable library paths.

\textbf{Format-semantic libraries (libsndfile, +9.27~pp).}
The smallest gain is on a library where the coverage bottleneck is not
handle dependency but rather semantic validity of byte payloads (audio
sample buffers, codec headers).  RNR fixes handle gaps but cannot
synthesize semantically valid byte content; this result correctly bounds
the scope of the repair mechanism.

The ablation isolates RNR as the specific component responsible for
converting a protobuf-schema harness that \emph{skips} dependents into
one that \emph{fixes} them, and quantifies the coverage cost of not
doing so.

%% =====================================================================
\section{Side-by-Side: libErator Fixed-Sequence Driver vs.\ Proto-libErator Super Harness}
\label{sec:eval_driver_comparison}
%% =====================================================================

Figure~\ref{fig:driver_comparison} (Appendix~\ref{app:thesis-alignment}) illustrates the fundamental difference between libErator's fixed-sequence driver model and Proto-libErator's Super Harness model.

In libErator, each driver is generated from a specific sequence of API calls inferred by the AFG traversal. The driver calls those APIs in a fixed order, with fixed argument roles. The fuzzer only varies argument \emph{values} — it cannot reorder the calls, insert new calls, or remove calls. If the optimal sequence for a specific code path requires a call ordering that no libErator driver happens to include, that path is unreachable.

In Proto-libErator, the protobuf schema encodes every API as an action variant. The \texttt{repeated Action actions} field gives the fuzzer full control over the sequence: it can reorder actions, repeat actions, omit actions, and vary action arguments — all within the bounds of the typed schema. The SCG-guided seeds and RNR guards bias the fuzzer toward valid orderings, but the fuzzer retains the freedom to explore.

The practical difference becomes most apparent on libraries with complex multi-step interaction patterns. For cjson, building a deeply nested JSON structure requires a sequence of \texttt{CreateObject}, \texttt{CreateArray}, multiple \texttt{AddItem} calls, and nested \texttt{AddItemToObject} calls in a specific interleaving. libErator's fixed-sequence drivers can capture this if the AFG traversal generates a driver covering the right path, but the search space of valid orderings grows exponentially with nesting depth. The Super Harness explores this space directly.

%% =====================================================================
\section{Constraint Graph Visualization: cjson Subgraph}
\label{sec:eval_cg_figure}
%% =====================================================================

Figure~\ref{fig:cjson_scg} (Appendix~\ref{app:thesis-alignment}) shows a representative subgraph of the cjson SCG, covering six core APIs and their constraint relationships. The full graph has 2,085 edges across 78 APIs; this subgraph is selected to illustrate the three edge types (producer--consumer hard, producer--consumer soft, and invalidation) in a compact and readable form.

The figure makes the SCG semantics concrete. Both \texttt{cJSON\_Parse} and \texttt{cJSON\_CreateObject} are confirmed producers (hard evidence, $\rho = 0.95$): they both return \texttt{struct cJSON*} and the Clang AST confirms the type. \texttt{cJSON\_GetObjectItem} is a weak producer ($\rho = 0.6$, soft): it returns \texttt{struct cJSON*} but only under the condition that the queried key exists, making its output unreliable as a handle source. The invalidation edges from \texttt{cJSON\_Delete} prevent any consumer (Add, Get, Print) from executing after a Delete without first producing a new handle — encoding the use-after-free prevention directly in the harness.


\subsection{Crash Signal and Triage Status}
\label{sec:eval_crashes}

\begin{table}[htb]
\centering
\caption{CASR-triaged crash clusters by bug class (Memory, Null/Ptr, Resource, Logic).}
\label{tab:crash_quality}
\resizebox{\columnwidth}{!}{%
\begin{tabular}{lrrrrrr}
\toprule
\textbf{Library} & \textbf{Raw} & \textbf{Clusters} & \textbf{Memory} & \textbf{Null/Ptr} & \textbf{Resource} & \textbf{Logic} \\
\midrule
cjson        &  36 &  29 &  8 & 21 &  0 & 0 \\
c-ares       &  38 &  35 &  8 & 29 &  1 & 0 \\
libhtp       &  71 &  67 &  2 & 64 &  5 & 0 \\
libpcap      &  81 &  75 &  9 & 66 &  4 & 2 \\
libsndfile   &   6 &   6 &  3 &  3 &  0 & 0 \\
libplist     & 179 & 155 & 79 & 90 &  2 & 8 \\
libtiff      & 182 & 153 & 21 &153 &  5 & 3 \\
pthreadpool  &   0 &   0 &  0 &  0 &  0 & 0 \\
\midrule
\textbf{Total} & \textbf{593} & \textbf{520} & \textbf{130} & \textbf{426} & \textbf{17} & \textbf{13} \\
\bottomrule
\end{tabular}
}
\end{table}

\subsection{Bug Validation and Provenance Classification}
\label{sec:eval_bug_validation}

The crash table above reports CASR clusters, which are useful for triage
but are not the unit of the paper's bug claim.  We treat a defect as
validated when it satisfies three conditions: (i) it reproduces under
AddressSanitizer on the target library revision used in the campaign;
(ii) the crashing path reaches library code rather than a harness-only
failure (confirmed by stack-frame inspection); and (iii) the triggering
input is minimized to a stable protobuf action sequence that can be
replayed outside the fuzzer loop.

Table~\ref{tab:bug_provenance} classifies all 36 validated crashes by
provenance, cross-referenced against the NVD CVE database, each
library's upstream bug tracker or security advisory list, and the
OSS-Fuzz issue tracker.  The provenance codes are: \textbf{N} = novel
(no matching public CVE or upstream issue found at classification time);
\textbf{K} = known-unpatched (an open upstream issue exists with no CVE
assigned and no patch shipped as of the campaign date); \textbf{V} =
needs version pinning (a related upstream fix landed recently; the exact
campaign binary version must be confirmed); and \textbf{S} = suspicious
(the crash site is a libc function such as \texttt{vsnprintf} or
\texttt{\_\_asan\_memcpy}, making harness-induced mis-attribution
plausible; full stack trace required for confirmation).

%% Bug provenance table -- generated from validated_bug_casebook_latest.md
%% Cross-referenced against NVD, OSS-Fuzz, and per-library upstream trackers.
%%
%% Provenance codes:
%%   N = Novel/Pending -- no matching CVE or public issue found; pending upstream disclosure
%%   K = Known-unpatched -- matches an open upstream issue with no CVE assigned, no patch shipped
%%   S = Suspicious -- crash site is a libc function (vsnprintf, __asan_memcpy); likely
%%       harness-induced artifact or mis-symbolized frame; requires stack-trace verification
%%   V = Verify -- a related fix landed in a recent upstream release; version pinning needed
%%       to confirm the campaign binary predates or postdates the fix
\begin{table}[htb]
\centering
\caption{Bug provenance classification for 36 validated crashes across 8 libraries.
  \textbf{N} = novel (no public CVE/issue); \textbf{K} = known-unpatched (open upstream issue,
  no CVE, no patch); \textbf{S} = suspicious (crash site in libc; possible harness-induced
  artifact); \textbf{V} = needs version verification.
  Bug class: \textbf{HBO-w} = heap-buffer-overflow (write); \textbf{UAF-w} = heap-use-after-free (write).}
\label{tab:bug_provenance}
\small
\setlength{\tabcolsep}{3pt}
\resizebox{\linewidth}{!}{%
\begin{tabular}{lllcl}
\toprule
\textbf{Library} & \textbf{Crash site (function)} & \textbf{Class} & \textbf{P} & \textbf{Notes} \\
\midrule
\multirow{6}{*}{c-ares}
  & \texttt{ares\_create\_query\_int}  & HBO-w & N & No CVE for this helper; CVE-2016-5180 was in public wrapper, fixed \\
  & \texttt{ares\_dns\_pton}           & HBO-w & N & New DNS API (v1.28+); no public report \\
  & \texttt{ares\_dns\_rr\_get\_keys}  & HBO-w & N & New DNS record accessor; no public report \\
  & \texttt{ares\_dns\_rr\_get\_opt\_byid} & HBO-w & N & New DNS record accessor; no public report \\
  & \texttt{ares\_expand\_name}        & HBO-w & N & No matching 2024--2026 CVE or commit \\
  & \texttt{ares\_version}             & HBO-w & S & Returns string constant + writes one int; heap write impossible by inspection \\
\midrule
\multirow{5}{*}{cjson}
  & \texttt{print\_string\_ptr}        & HBO-w & K & GitHub issue \#923 (Jan 2025), open, no CVE, no patch in v1.7.18/v1.7.19 \\
  & \texttt{print\_array}              & HBO-w & K & Same root cause as \#923 (print subsystem buffer sizing) \\
  & \texttt{print\_number}             & HBO-w & V & v1.7.19 added number-buffer fix; confirm campaign binary version \\
  & \texttt{print\_object}             & HBO-w & K & Same root cause as \#923 \\
  & \texttt{cJSON\_ReplaceItemViaPointer} & UAF-w & N & UAF write via \texttt{cJSON\_ReplaceItemInArray}; no CVE or fix commit \\
\midrule
\multirow{10}{*}{libdwarf}
  & \texttt{\_dwarf\_get\_fde\_list\_internal}           & HBO-w & N & Not in official bug DB (prevanders.net/dwarfbug.html) \\
  & \texttt{dwarf\_get\_loclist\_offset\_index\_value}   & HBO-w & N & No match in NVD or libdwarf bugxml \\
  & \texttt{dwarf\_get\_rnglist\_offset\_index\_value}   & HBO-w & N & No match in NVD or libdwarf bugxml \\
  & \texttt{dwarf\_get\_section\_info\_by\_index\_a}~(×2) & HBO-w & N & Two distinct crash PCs; no public report \\
  & \texttt{dwarf\_get\_section\_info\_by\_name\_a}~(×4) & HBO-w & N & Four distinct crash PCs; no public report \\
  & \texttt{dwarf\_globals\_by\_type}                    & HBO-w & N & DW202508-001 fixed \texttt{dwarf\_globals.c} in Aug 2025; distinct entry point \\
  & \texttt{dwarf\_hasattr}                              & HBO-w & N & No match in NVD or libdwarf bugxml \\
  & \texttt{dwarf\_offset\_list}                         & HBO-w & N & No match in NVD or libdwarf bugxml \\
  & \texttt{dwarf\_uncompress\_integer\_block\_a}        & HBO-w & N & No match in NVD or libdwarf bugxml \\
\midrule
libhtp
  & \texttt{bstr\_alloc} (bstr.c:48)  & HBO-w & N & No CVE; published advisories cover resource exhaustion, not allocator OOB \\
\midrule
\multirow{3}{*}{libpcap}
  & \texttt{pcap\_lookupnet}           & HBO-w & N & Not in 2024--2026 CVE list (CVE-2023-7256, CVE-2024-8006, CVE-2025-11961 unrelated) \\
  & \texttt{pcapint\_parsesrcstr\_ex}  & HBO-w & N & Internal API; no public report \\
  & \texttt{vsnprintf} (libpcap ctx)   & HBO-w & S & libc function; errbuf sizing or harness scaffolding likely source \\
\midrule
\multirow{3}{*}{libplist}
  & \texttt{plist\_from\_memory}       & HBO-w & N & Issue \#244 was a READ OOB (fixed 2023); write OOB is distinct \\
  & \texttt{node\_attach}              & UAF-w & N & No public CVE; issue \#276 named different functions \\
  & \texttt{node\_insert}              & UAF-w & N & No public CVE; sibling to \texttt{node\_attach} \\
\midrule
libsndfile
  & \texttt{vsnprintf} (sndfile ctx)   & HBO-w & S & libc function; crash in internal log-format buffer; harness involvement likely \\
\midrule
\multirow{4}{*}{libtiff}
  & \texttt{LogLuv24toXYZ}             & HBO-w & N & 2016-era LUV codec fix unrelated; no 2024--2026 CVE for this function \\
  & \texttt{LogLuv32toXYZ}             & HBO-w & N & Same as above \\
  & \texttt{TIFFYCbCrToRGBInit}        & HBO-w & N & 2016/2017 int overflow fixed in v4.0.8; write OOB variant has no matching CVE \\
  & \texttt{\_TIFFGetStrileOffset\ldots} & HBO-w & N & tif\_dirread.c; no 2024--2026 CVE for this accessor \\
\midrule
\multicolumn{2}{l}{\textbf{Total}} & & & \textbf{29~N, 3~K, 1~V, 3~S} \\
\bottomrule
\end{tabular}
}
\end{table}


The classification yields 29~novel findings (N), 3 known-but-unpatched
findings in cjson's print subsystem (K, related to GitHub issue~\#923,
opened January 2025, no CVE assigned, no patch in any release as of the
campaign date), 1 finding requiring version confirmation (V), and
3 suspicious crash sites that may be harness-induced artifacts (S).

\textbf{Novel findings by library.}
The 13 libdwarf findings span seven distinct API entry points covering
frame-description lists, location/range index values, section enumeration,
global name tables, attribute lookup, offset lists, and integer-block
decompression; none appear in the official libdwarf vulnerability
database (\textit{prevanders.net/dwarfbug.html}) or any NVD entry.
The 5 c-ares findings target the new DNS record API (v1.28+) introduced
in the library's 2023 rewrite; functions \texttt{ares\_dns\_pton},
\texttt{ares\_dns\_rr\_get\_keys}, and \texttt{ares\_dns\_rr\_get\_opt\_byid}
have no matching security advisories.
The libplist pair of use-after-free writes in \texttt{node\_attach}
and \texttt{node\_insert} are distinct from the previously documented
UAF in \texttt{plist\_free} (issue~\#276) and from the 2023 read-OOB
in \texttt{plist\_from\_memory} (issue~\#244).
The 4 libtiff findings are in the LogLuv codec, YCbCr conversion
initialization, and TIFF directory reader; none match any CVE from 2024
onward (CVE-2025-8177, CVE-2025-8851, CVE-2025-9900, CVE-2026-4775 all
target different code paths).

\textbf{Suspicious crash sites excluded from novelty claim.}
Three of the 36 crashes land in \texttt{vsnprintf} (libpcap and
libsndfile) or are attributed to \texttt{\_\_asan\_memcpy} (cjson
Case~7).  These warrant caution: \texttt{vsnprintf} is a libc function
and a crash attributed to it likely reflects an imprecise symbolization
of the real fault site; it may also be harness-induced if the harness
passes an incorrectly sized error buffer.  The \texttt{ares\_version}
crash (Case~6) is also excluded from the novelty claim because the
function returns only a string constant and writes one integer
constant --- a heap-buffer-overflow write there is impossible by source
inspection, and the crash is almost certainly a harness artifact.
These four cases are noted as requiring additional stack-trace inspection
before upstream reporting.

\subsection{Claim Boundaries and Ongoing Tasks}
\label{sec:eval_remaining}

The current evidence supports three class-dependent claims: (i) a static head-start on stateful/constraint-dense libraries via SCG-guided seeds, (ii) broader sequence exploration via typed super-harness dispatch, and (iii) the RNR ablation (Table~\ref{tab:ablation_12libs}) directly isolates the coverage cost of skipping dependents rather than repairing them---SCG and LPM are held constant across both ablation conditions, so the gains are attributable to the repair mechanism alone. Remaining tasks (equal-budget baseline reruns, broader causal-edge auditing, and completion of the public bug-provenance table) are tracked as follow-up work rather than incorporated as completed claims.
